% !TeX program = lualatex
% !TeX encoding = utf8
% !TeX spellcheck = uk_UA
% !TeX root = main.tex

%% --------------------------------------------------------
%% Додаток: перехід між гауссовою системою (СГС) і СІ, константи
%% (за матеріалами Notes/SI2SGS.tex; таблиця формул плазми додана)
%% --------------------------------------------------------

\section{Перехід між СГС і~СІ. Константи}\label{sec:SGStoSI}\hypertarget{app-sgs-si}{}

У~цьому посібнику використано гауссову систему (СГС), а~значна частина сучасної навчальної й~прикладної
літератури, зокрема з~термоядерних установок, спирається на~СІ. Перехід між системами~--- це не~просте
масштабування, а~зміна вигляду рівнянь: у~СІ з'являються сталі $\varepsilon_0$ і~$\mu_0$, а~в~гауссовій системі
їх нема, і~відповідні коефіцієнти стоять у~формулах як множники $4\pi$ та~$c$. Тому формули з~СГС не~можна
просто «перекласти» заміною символів. Нижче наведено числові коефіцієнти переведення одиниць, вигляд основних
формул плазми в~обох системах і~найважливіші константи.

\subsection*{Числові коефіцієнти переведення}

Остання колонка таблиці~\ref{tab:SI-to-CGS} дає, у~скільки разів числове значення величини в~СГС більше,
ніж у~СІ: наприклад, $1~\text{Тл}=10^{4}$~Гс, $1~\text{Дж}=10^{7}$~ерг. Швидкість світла в~цих коефіцієнтах
округлено до $3\times10^{10}$~см/с.

\begin{longtblr}[
		caption={Числові коефіцієнти переведення між СІ та~СГС},
		label={tab:SI-to-CGS},
	]{colspec={X[2.4,l,m]X[1.2,l,m]X[1.5,l,m]X[1.2,c,m]}, rowhead=1}
	Величина & Система СІ & Система СГС & СГС/СІ \\
	\SetCell[c=4]{l,m,bg=themecolorlight!16,fg=themecolordark,font=\bfseries\sffamily\small} Основні та механічні одиниці & & & \\
	Довжина & метр, м & сантиметр, см & $10^{2}$ \\
	Маса & кілограм, кг & грам, г & $10^{3}$ \\
	Час & секунда, с & секунда, с & $1$ \\
	Сила & ньютон, Н & дина & $10^{5}$ \\
	Робота, енергія & джоуль, Дж & ерг & $10^{7}$ \\
	Тиск & паскаль, Па & дин/см$^2$ & $10$ \\
	\SetCell[c=4]{l,m,bg=themecolorlight!16,fg=themecolordark,font=\bfseries\sffamily\small} Електричні одиниці & & & \\
	Електричний заряд & кулон, Кл & франклін, Фр (СГСЕ) & $3\times10^{9}$ \\
	Густина заряду & Кл/м$^3$ & Фр/см$^3$ & $3\times10^{3}$ \\
	Сила струму & ампер, А & Фр/с (статампер) & $3\times10^{9}$ \\
	Густина струму & А/м$^2$ & Фр/(с$\cdot$см$^2$) & $3\times10^{5}$ \\
	Дипольний момент & Кл$\cdot$м & Фр$\cdot$см & $3\times10^{11}$ \\
	Напруга, потенціал & вольт, В & статвольт, статВ & $\tfrac13\times10^{-2}$ \\
	Напруженість поля $\vect E$ & В/м & статВ/см & $\tfrac13\times10^{-4}$ \\
	Індукція поля $\vect D$ & Кл/м$^2$ & Фр/см$^2$ & $12\pi\times10^{5}$ \\
	Опір & ом, Ом & с/см & $\tfrac19\times10^{-11}$ \\
	Ємність & фарад, Ф & см & $9\times10^{11}$ \\
	\SetCell[c=4]{l,m,bg=themecolorlight!16,fg=themecolordark,font=\bfseries\sffamily\small} Магнітні одиниці & & & \\
	Напруженість поля $\vect H$ & А/м & ерстед, Е & $4\pi\times10^{-3}$ \\
	Магнітна індукція $\vect B$ & тесла, Тл & гаус, Гс & $10^{4}$ \\
	Магнітний момент & А$\cdot$м$^2$ & ерг/Гс & $10^{3}$ \\
	Магнітний потік & вебер, Вб & максвел, Мкс & $10^{8}$ \\
	Індуктивність & генрі, Гн & см & $10^{9}$ \\
\end{longtblr}

\subsection*{Основні формули фізики плазми в~СГС і~СІ}

Таблиця~\ref{tab:plasma-formulas-units} показує, як виглядають формули цього посібника в~обох системах.
Температура скрізь в~енергетичних одиницях; заряд іона для простоти $Ze=e$.

\begin{longtblr}[
		caption={Основні формули фізики плазми в~гауссовій системі та~в~СІ},
		label={tab:plasma-formulas-units},
	]{colspec={X[1.7,l,m]X[c,m]X[c,m]}, rowhead=1}
	Величина & СГС & СІ \\
	Сила Лоренца & $q\bigl(\vect E+\vect v\times\vect B/c\bigr)$ & $q\bigl(\vect E+\vect v\times\vect B\bigr)$ \\
	Плазмова частота $\omega_{pe}^2$ & $\dfrac{4\pi n e^2}{m_e}$ & $\dfrac{n e^2}{\varepsilon_0 m_e}$ \\
	Дебаївська довжина $\lambda_{De}^2$ & $\dfrac{T_e}{4\pi n e^2}$ & $\dfrac{\varepsilon_0 T_e}{n e^2}$ \\
	Циклотронна частота $\omega_{c}$ & $\dfrac{eB}{mc}$ & $\dfrac{eB}{m}$ \\
	Ларморівський радіус $r_L$ & $\dfrac{m v_\perp c}{eB}$ & $\dfrac{m v_\perp}{eB}$ \\
	Критична концентрація $n_{\text{кр}}$ & $\dfrac{m_e\omega^2}{4\pi e^2}$ & $\dfrac{\varepsilon_0 m_e\omega^2}{e^2}$ \\
	Енергія взаємодії пари на~відстані $r$ & $\dfrac{e^2}{r}$ & $\dfrac{e^2}{4\pi\varepsilon_0 r}$ \\
	Радіус розсіяння на~$90^\circ$ $b_{90}$ & $\dfrac{e^2}{m_e v^2}$ & $\dfrac{e^2}{4\pi\varepsilon_0 m_e v^2}$ \\
	Частота зіткнень $\nu_{ei}$ & $\dfrac{4\sqrt{2\pi}}{3}\,\dfrac{n e^4\ln\Lambda}{m_e^{1/2}T_e^{3/2}}$ & $\dfrac{\sqrt{2\pi}}{12\pi^2}\,\dfrac{n e^4\ln\Lambda}{\varepsilon_0^2 m_e^{1/2}T_e^{3/2}}$ \\
	Поле Драйсера $E_D$ & $\dfrac{4\pi e^3 n\ln\Lambda}{T_e}$ & $\dfrac{e^3 n\ln\Lambda}{4\pi\varepsilon_0^2 T_e}$ \\
	Магнітний тиск & $\dfrac{B^2}{8\pi}$ & $\dfrac{B^2}{2\mu_0}$ \\
	Параметр $\beta$ & $\dfrac{8\pi p}{B^2}$ & $\dfrac{2\mu_0 p}{B^2}$ \\
	Альвенівська швидкість $v_A$ & $\dfrac{B}{\sqrt{4\pi\rho}}$ & $\dfrac{B}{\sqrt{\mu_0\rho}}$ \\
	Магнітне число Рейнольдса $R_m$ & $\dfrac{4\pi\sigma L v}{c^2}$ & $\mu_0\sigma L v$ \\
	Коефіцієнт дифузії поля & $\dfrac{c^2}{4\pi\sigma}$ & $\dfrac{1}{\mu_0\sigma}$ \\
\end{longtblr}

Безрозмірні комбінації ($\beta$, $R_m$, $N_D$, $\Gamma$, відношення частот) в~обох системах однакові~---
різняться лише записи через розмірні величини. Тому порівняння з~літературою зручно починати саме з~безрозмірних параметрів.

\subsection*{Фізичні константи}

Значення найуживаніших констант у~гауссовій системі зібрано в~таблиці~\ref{tab:phys-constants}
(за~рекомендаціями CODATA).

\begin{longtblr}[
		caption={Деякі фізичні константи},
		label={tab:phys-constants},
	]{colspec={X[1.6,l,m]X[0.7,c,m]X[2,l,m]}, rowhead=1}
	Константа & Символ & Значення \\
	\SetCell[c=3]{l,m,bg=themecolorlight!16,fg=themecolordark,font=\bfseries\sffamily\small} Фундаментальні сталі & & \\
	Швидкість світла у~вакуумі & $c$ & $2{,}997\,924\,58\times10^{10}$~см/с \\
	Гравітаційна стала & $G$ & $6{,}674\times10^{-8}$~см$^3$/(г$\cdot$с$^2$) \\
	Стала Планка & $\hbar$ & $1{,}054\,572\times10^{-27}$~ерг$\cdot$с \\
	Елементарний заряд & $e$ & $4{,}803\,204\times10^{-10}$~Фр \\
	\SetCell[c=3]{l,m,bg=themecolorlight!16,fg=themecolordark,font=\bfseries\sffamily\small} Елементарні частинки & & \\
	Маса електрона & $m_e$ & $9{,}109\,384\times10^{-28}$~г \\
	Енергія спокою електрона & $m_ec^2$ & $0{,}511$~МеВ \\
	Маса протона & $m_p$ & $1{,}672\,622\times10^{-24}$~г \\
	Енергія спокою протона & $m_pc^2$ & $938{,}27$~МеВ \\
	Відношення мас & $m_p/m_e$ & $1836{,}15$ \\
	\SetCell[c=3]{l,m,bg=themecolorlight!16,fg=themecolordark,font=\bfseries\sffamily\small} Атомні й~термодинамічні сталі & & \\
	Борівський радіус & $a_0$ & $5{,}291\,772\times10^{-9}$~см \\
	Магнетон Бора & $\mu_B$ & $9{,}274\,010\times10^{-21}$~ерг/Гс \\
	Стала Больцмана & $k_{\mathrm B}$ & $1{,}380\,649\times10^{-16}$~ерг/К;\quad $8{,}617\,333\times10^{-5}$~еВ/К \\
	Універсальна газова стала & $R$ & $8{,}314\times10^{7}$~ерг/(К$\cdot$моль) \\
	Число Авогадро & $N_A$ & $6{,}022\times10^{23}$~моль$^{-1}$ \\
	Атомна одиниця маси & $u$ & $1{,}660\,539\times10^{-24}$~г \\
\end{longtblr}

\subsection*{Енергетичні еквіваленти}

Електронвольт~--- основна енергетична одиниця фізики плазми, тому корисно мати під рукою його еквіваленти
(таблиця~\ref{tab:eV-equivalents}).

\begin{longtblr}[
		caption={Еквіваленти електронвольта},
		label={tab:eV-equivalents},
	]{colspec={X[1.4,l,m]X[2,c,m]X[2,c,m]}, rowhead=1}
	Величина & $1$~еВ відповідає & Обернено \\
	Енергія & $1{,}602\times10^{-12}$~ерг $=1{,}602\times10^{-19}$~Дж & $1$~ерг $=6{,}242\times10^{11}$~еВ \\
	Температура & $11\,605$~К & $1$~К $=8{,}617\times10^{-5}$~еВ \\
	Частота & $2{,}418\times10^{14}$~Гц & $1$~Гц $=4{,}136\times10^{-15}$~еВ \\
	Довжина хвилі фотона & $1{,}240$~мкм & $1$~мкм $\leftrightarrow1{,}240$~еВ \\
	Маса & $1$~еВ$/c^2=1{,}783\times10^{-33}$~г & $1$~г $=5{,}610\times10^{32}$~еВ$/c^2$ \\
\end{longtblr}

\clearpage
